\subsection{Transport of the two compact shifts to the actual orbital norm}
\label{actual-wigner-orbital-transport}

Retain the original full-group cofactor construction from
\emph{A regular coefficient for the inverse terminal operator}
(the construction denoted \texttt{rl-section}, also used in
the actual top-lift proof), the place data, its quotient measures,
and the two compact maps from the compact Wigner theorem.
Write $s=r+1$ and let $a\ge1$ be their
chain length (in the equal-type application, $a=30s^2+6$).
Identify $W'_+=W\cap(z')^\perp$ with $\mathbb C^k$ by a fixed
orthonormal basis.  Let the original right compact representation
be $\pi'$ on $W_0$, and the shifted representations be $\pi_+$
on $W_+$ and $\pi_-$ on $W_-$.  The compact theorem supplies
polynomial maps
\[
 A_\pm:W'_+\longrightarrow\operatorname{Hom}(W_\pm,W_0)
\]
of total degree at most $a$ in the real coordinates (equivalently,
in the complex coordinates and their conjugates), with
\begin{align}
 A_\pm(h\omega)&=\pi'(h)A_\pm(\omega)\pi_\pm(h)^{-1},
                                                        \label{wigner-transport-cov}\\
 \|A_\pm(\omega)\|&\le\|\omega\|^a,\label{wigner-transport-upper}\\
 \|\omega\|^{2a}\|v\|^2
 &\le\kappa^2\bigl(\|A_+(\omega)^*v\|^2
                         +\|A_-(\omega)^*v\|^2\bigr)
 \quad(v\in W_0).                                      \label{wigner-transport-lower}
\end{align}
Here $h\in K_{z'}\simeq U(k)$, and the compact theorem gives
$\kappa=(3b)^a$.  These maps may be transported by unitary
intertwiners from a highest-weight model to the actual compact
representations; every displayed norm and covariance is preserved.

Let $\Omega(g)\in W'_+$ be the actual cofactor vector and define
\[
              \mathcal A_\pm(g)=A_\pm(\Omega(g)).
\]
The proved cofactor identities are
\begin{align}
 \|\Omega(g)\|^2&=\mathcal G(g)/T',\label{wigner-transport-gram}\\
 \Omega(k_L^{-1}n_L^{-1}g n_Rk_R)&=k_R^{-1}\Omega(g),
                                                        \label{wigner-transport-omega}
\end{align}
for the original $k_L\in K_z$, $k_R\in K_{z'}$ and both
unipotent radicals.  No intermediate minor is assumed nonzero
in these identities.  Indeed the triangular action of the two
unipotent groups on the indicated frames gives
$Z(g')=A^*Z(g)B$ with $A,B$ upper triangular unipotent;
$\operatorname{adj}(Z)Z=\det(Z)I$ and the corresponding
adjugate transformation show that the cofactor vector transforms
by $k_R^{-1}n_R^{-1}$ after the left action cancels.
The projection to $W'_+$ removes the $U$-radical part of
$n_R^{-1}$, leaving exactly
\eqref{wigner-transport-omega}.  These polynomial identities
extend across vanishing minors.  Since each coordinate of $\Omega$ is a
polynomial of degree $s$ in the conjugate matrix entries,
every matrix entry of $\mathcal A_\pm$ is a polynomial of
degree at most $as$ on the full group.  In particular no inverse
Gram factor occurs in these shift maps.

For a coefficient $F(g):W_0\to W_p$, set
$F_\pm(g)=F(g)\mathcal A_\pm(g)$, and use the shifted right
type $\pi_\pm$ in the corresponding original orbital integral.

\begin{proposition}\label{wigner-transport-orbital}
At every valid original representative $g_\gamma$ with
$\mathcal G(g_\gamma)>0$, the orbital maps obey
\begin{equation}\label{wigner-transport-orbital-factor}
 \operatorname{Orb}_\pm(F_\pm;g_\gamma)
       =\operatorname{Orb}_0(F;g_\gamma)
                         A_\pm(\Omega(g_\gamma)).
\end{equation}
This equality has the original quotient and compact Haar
normalizations on both sides.  Let $E_0(F;\gamma)$ and
$E_\pm(F_\pm;\gamma)$ denote the original empty-word point
energies in these three compact types.  Then
\begin{align}
 (\mathcal G(\gamma)/T')^a E_0(F;\gamma)
 &\le\kappa^2\{E_+(F_+;\gamma)+E_-(F_-;\gamma)\},
                                                        \label{wigner-transport-point-lower}\\
 E_\pm(F_\pm;\gamma)
 &\le(\mathcal G(\gamma)/T')^a E_0(F;\gamma).
                                                        \label{wigner-transport-point-upper}
\end{align}
\end{proposition}

\begin{proof}
Write $g'=k_L^{-1}n_L^{-1}g_\gamma n_Rk_R$.  Equations
\eqref{wigner-transport-cov} and
\eqref{wigner-transport-omega} give
\[
 \mathcal A_\pm(g')
   =\pi'(k_R)^{-1}A_\pm(\Omega(g_\gamma))\pi_\pm(k_R).
\]
The integrand for the shifted orbital integral, before its common
unit-modulus character, is therefore
\begin{align*}
 \pi(k_L)F(g')\mathcal A_\pm(g')\pi_\pm(k_R)^{-1}
 &=\pi(k_L)F(g')\pi'(k_R)^{-1}
                                    A_\pm(\Omega(g_\gamma)).
\end{align*}
The final matrix is independent of every integration variable.
It factors out on the right, proving
\eqref{wigner-transport-orbital-factor}.  The character, both
unipotent integrations, compact averages, and discarded-stabilizer
quotient cut are unchanged.  For compact coefficients this is
ordinary integration of the exact integrand identity.  The same
identity applies to the prescribed finite-decay coefficients by
the established absolute quotient-integrability bounds at each
fixed valid representative, or by their proved compact exhaustion.

Put $B=\operatorname{Orb}_0(F;g_\gamma)$.  Applying
\eqref{wigner-transport-lower} to $B^*e$ for an orthonormal
basis $e$ of $W_p$ and summing gives
\[
 \|\Omega\|^{2a}\|B\|_{\rm HS}^2
 \le\kappa^2(\|BA_+(\Omega)\|_{\rm HS}^2
                       +\|BA_-(\Omega)\|_{\rm HS}^2).
\]
Likewise \eqref{wigner-transport-upper} gives
$\|BA_\pm(\Omega)\|_{\rm HS}^2
 \le\|\Omega\|^{2a}\|B\|_{\rm HS}^2$.
The original point-density factor depends on $k,r,T,T'$ and
$\gamma$, all unchanged by the compact-type shift.  Multiply
the Hilbert--Schmidt inequalities by that common nonnegative
factor and use \eqref{wigner-transport-gram}.  This proves
both energy estimates without a representation-dimension factor.
\end{proof}

\begin{corollary}\label{wigner-transport-moment}
For every integer $p\ge0$, whenever the indicated finite-decay
norms are defined,
\begin{equation}\label{wigner-transport-norm-upper}
 jN_\pm^2(\mathcal G^pF\mathcal A_\pm)
       \le (T')^{-a}jN_0^2(D_s^{\,2p+a}F).
\end{equation}
The squared norm on the right is the original leading-minor
Jacquet norm in the original compact type.
\end{corollary}

\begin{proof}
The polynomial $\mathcal G$ is invariant under both original
compact and unipotent actions.  Multiplication by
$\mathcal G^p$ therefore multiplies the orbital map at $\gamma$
by $\mathcal G(\gamma)^p$.  On the positive domain,
\[
 0<\mathcal G=|D_s|^2-TT'|D_{s-1}|^2\le|D_s|^2.
\]
Multiplying \eqref{wigner-transport-point-upper} by
$\mathcal G^{2p}$ gives
\[
 E_\pm(\mathcal G^pF\mathcal A_\pm;\gamma)
 \le (T')^{-a}|D_s(\gamma)|^{2(2p+a)} E_0(F;\gamma).
\]
The complex polynomial $D_s$ is itself invariant under the same
actions, so the right side is exactly
$(T')^{-a}E_0(D_s^{2p+a}F;\gamma)$.  Integrate over the
original big-cell norm variables with their original measures.
The point-density is the same in both compact types.  Tonelli
proves the inequality even before finiteness; the stated moment
bound supplies finiteness in the application.
\end{proof}

\paragraph{Cancellation in the equal-type estimate.}
Suppose each shifted type has the established different-type
pointwise inequality with the same numerical factor $C$:
\[
 E_\pm(F_\pm;\gamma)
       \le C\,\mathcal G(\gamma)^{-2p}
                        jN_\pm^2(\mathcal G^pF_\pm).
\]
Combining the proposition and corollary gives the exact bound
\begin{equation}\label{wigner-transport-combined}
 E_0(F;\gamma)
 \le 2C\kappa^2\mathcal G(\gamma)^{-2p-a}
                              jN_0^2(D_s^{2p+a}F).
\end{equation}
The $T'^a$ from the pointwise lower bound cancels the
$(T')^{-a}$ from the norm comparison.  With $p=6s^2$ and
$a=30s^2+6$, the inverse-Gram exponent is
$2p+a=42s^2+6\le48s^2$ and the polynomial degree required
by the spectral decomposition is
$2sp+as=42s^3+6s\le48s^3$.
Thus both retain the required quadratic pole order and cubic
spectral reserve.  The existing leading-minor moment theorem
and the compact bound $\kappa=(3b)^a$ supply the remaining
uniform factor in the original unshifted norm.
